%%
%% This is file `sample-acmcp.tex',
%% generated with the docstrip utility.
%%
%% The original source files were:
%%
%% samples.dtx  (with options: `all,journal,acmcp')
%%
%% IMPORTANT NOTICE:
%%
%% For the copyright see the source file.
%%
%% Any modified versions of this file must be renamed
%% with new filenames distinct from sample-acmcp.tex.
%%
%% For distribution of the original source see the terms
%% for copying and modification in the file samples.dtx.
%%
%% This generated file may be distributed as long as the
%% original source files, as listed above, are part of the
%% same distribution. (The sources need not necessarily be
%% in the same archive or directory.)
%%
%%
%% Commands for TeXCount
%TC:macro \cite [option:text,text]
%TC:macro \citep [option:text,text]
%TC:macro \citet [option:text,text]
%TC:envir table 0 1
%TC:envir table* 0 1
%TC:envir tabular [ignore] word
%TC:envir displaymath 0 word
%TC:envir math 0 word
%TC:envir comment 0 0
%%
%% The first command in your LaTeX source must be the \documentclass
%% command.
%%
%% For submission and review of your manuscript please change the
%% command to \documentclass[manuscript, screen, review]{acmart}.
%%
%% When submitting camera ready or to TAPS, please change the command
%% to \documentclass[sigconf]{acmart} or whichever template is required
%% for your publication.
%%
%%
\documentclass[sigconf,review,nonacm]{acmart}
%%
%% \BibTeX command to typeset BibTeX logo in the docs
\AtBeginDocument{%
  \providecommand\BibTeX{{%
    Bib\TeX}}}

%% Rights management information.
\setcopyright{none}
\copyrightyear{2026}
\acmYear{2026}
\acmDOI{XXXXXXX.XXXXXXX}
\acmConference[WSDM '26]{Proceedings of the 19th ACM International Conference on Web Search and Data Mining (WSDM)}{February 22--26, 2026}{Boise, ID}

% acmart already loads hyperref; avoid reloading.
\usepackage{fontawesome5}
\usepackage{float} % for [H] exact placement

\begin{document}

\title{LLM-to-Map: Conversational Co-Simulation of Power and Transportation Networks via LLM-Driven Map Control}

\author{Marouane Benbrahim}
\email{mb4194@msstate.edu}
\affiliation{%
  \institution{Mississippi State University}
  \state{Mississippi}
  \country{USA}
}

\author{Zonghan Zhang}
\email{zz239@msstate.edu}
\affiliation{%
  \institution{Mississippi State University}
  \state{Mississippi}
  \country{USA}
}

\author{Zhiqian Chen}
\email{zchen@cse.msstate.edu}
\affiliation{%
  \institution{Mississippi State University}
  \state{Mississippi}
  \country{USA}
}

\renewcommand{\shortauthors}{Benbrahim et al.}
\acmArticleType{Review}
% \keywords{Traffic-Electricity Co-Simulation, LLM}

\begin{abstract}
We present \textbf{LLM-to-Map}, a conversational interface for controlling and exploring a co-simulation of urban infrastructures. The system integrates SUMO for microscopic traffic modeling and PyPSA for electrical grid simulation, synchronized through an event-driven coordination engine that captures cascading effects between domains. A novel temporal dynamics system enables continuous time progression (e.g., 08:00 $\rightarrow$ 08:01 $\rightarrow$ 08:02) with intelligent vehicle quota management that automatically maintains realistic traffic density as vehicles reach destinations. An LLM agent powered by GPT-4 provides natural language interaction, translating user queries into simulator actions and visualization commands. LLM-to-Map supports fault-tolerant geospatial search, infrastructure control (e.g., triggering substation failures, activating vehicle-to-grid responses), and interactive map navigation. The Mapbox visualization layer renders disruptions and recovery processes in real-time. We showcase three scenarios: (1) grid outage cascading into traffic disruptions, (2) EV charging stress on substations, and (3) V2G-enabled resilience. By combining adaptive co-simulation with conversational control, LLM-to-Map lowers the barrier to studying cross-network resilience. Check our demo at \href{https://github.com/XGraph-Team/Map-LLM}{\faGithub\ Github.com/XGraph-Team/Map-LLM}.
\end{abstract}

\maketitle

\section{Introduction}
Urban infrastructures rarely operate in isolation. Transportation networks depend on electricity for traffic signals and EV charging, while power grids rely on mobility patterns for operation and recovery. Disruptions in one domain cascade into another: substation failures disable intersections, and traffic congestion delays grid repair. Domain-specific simulators like SUMO~\cite{lopez2018sumo} and PyPSA~\cite{brown2018pypsa} excel within their domains, and prior co-simulation frameworks such as GEMINI~\cite{panossian2023architecture} and V2Sim~\cite{qian2024v2sim} have shown promise in coupling networks. However, these platforms require technical expertise and restrict interaction to fixed parameters or scripted runs.

Large language models (LLMs) enable natural, conversational interfaces for complex systems. LLM agents have been explored in smart city decision support~\cite{kalyuzhnaya2025llm} and urban planning~\cite{han2025large}, with surveys highlighting their potential in multi-agent control~\cite{guo2024large}.

We introduce \textbf{LLM-to-Map}, combining co-simulation of power and traffic networks with an LLM-driven conversational interface and continuous temporal dynamics. Users issue commands such as "turn off Penn Station substation" or "set time to 8" and observe real-time impacts across both systems. The progressive time system initiates temporal flow where simulation time advances continuously (08:00 $\rightarrow$ 08:01...), while intelligent vehicle quota management maintains realistic traffic patterns by automatically spawning replacements as vehicles complete trips. This transforms static scenario execution into a living urban simulation supporting long-term resilience studies. The demo showcases grid-traffic cascades, EV charging stress, and V2G stabilization, offering a new paradigm for AI-assisted urban infrastructure analysis.

\section{System Overview}
The \textbf{LLM-to-Map} framework consists of three core layers: (1) domain simulators, (2) coordination engine with temporal dynamics, and (3) interaction and visualization layer (Figure~\ref{fig:structure}).

\noindent\textbf{Domain Simulators.}
We integrate SUMO for microscopic transport modeling~\cite{lopez2018sumo} and PyPSA for electrical network analysis~\cite{brown2018pypsa}. SUMO models central Manhattan (34th–59th Street) with hundreds of edges, intersections, and actuated traffic signals. Electric vehicles are explicitly represented with battery state-of-charge tracking for charging events. PyPSA represents eight major distribution substations (600–1000 MVA) connected to transmission nodes, distribution transformers, and downstream loads, accounting for both baseline demand and dynamic EV charging loads with power flow solved every five seconds.

\noindent\textbf{Temporal Dynamics and Resource Management.}
Unlike static scenario runners, LLM-to-Map implements a progressive time system where setting a time (e.g., "set time to 8") initiates continuous temporal flow (08:00 $\rightarrow$ 08:01 $\rightarrow$ 08:02...). The system maintains nine time-based traffic patterns (late night, early morning, morning rush, midday, afternoon, evening rush, evening, night, late evening), each defining vehicle count quotas (10–100 vehicles) and EV penetration rates (50–75\%). An intelligent quota manager monitors vehicle counts every 5 seconds, automatically spawning replacements when vehicles reach destinations, ensuring realistic traffic density throughout simulation. For example, during morning rush (7–9 AM), the system maintains 85–100 active vehicles; if count drops to 80, it spawns 5 more vehicles with appropriate EV ratios. This transforms the platform from a scenario player into a living, self-sustaining urban simulation, enabling long-term resilience studies and emergent behavior analysis without manual intervention.

\noindent\textbf{Event-Driven Coordination.}
A Python coordination engine synchronizes SUMO and PyPSA through a shared clock. Events propagate across domains: substation failures disable traffic lights forcing cautious vehicle flow, EV charging sessions aggregate into substation loads, and line overloads trigger alerts or shedding. This event-driven design preserves causality and supports cascading effects in real time. The backend uses Python Flask with SUMO and PyPSA running as separate processes. Natural language interaction is powered by GPT-4~\cite{achiam2023gpt} through the OpenAI API, orchestrated with LangChain. User queries are preprocessed with fuzzy-string matching (Levenshtein distance $\leq$2 edits) for typo-tolerant location retrieval.

\noindent\textbf{LLM-to-Map Interaction Layer.}
A GPT-4 + LangChain agent processes natural language instructions and maps them to infrastructure or visualization commands. Example interactions include "turn off Penn Station substation" (grid node failure cascades), "show me Times Square" (camera pans and zooms in Mapbox), and "activate vehicle-to-grid rescue" (EVs discharge energy). Intents are categorized into \emph{location search}, \emph{infrastructure control}, \emph{data query}, and \emph{navigation}.

\begin{table}[t]
\centering
\caption{System Comparison: Co-Simulation + Interaction Capabilities}
\small
\begin{tabular}{lccc}
\toprule
\textbf{Feature} & \textbf{LLM-to-Map} & \textbf{GEMINI} & \textbf{V2Sim} \\
\midrule
Traffic Simulator & SUMO & Custom & VISSIM \\
Grid Simulator & PyPSA & GridLAB-D & Custom \\
Natural Language & GPT-4 & No & No \\
Progressive Time & Yes & No & No \\
Auto Vehicle Quota & Yes & No & No \\
Typo Tolerance & Yes & No & No \\
Live 2D Viz & Mapbox & No & Limited \\
\bottomrule
\end{tabular}
\end{table}

\noindent\textbf{Visualization and Playback.}
The Mapbox interface renders the Manhattan testbed in 2D, overlaying grid nodes, traffic signals, and EV stations. Disruptions and recovery actions are animated with flashing outages, EV discharges, and smooth cinematic camera transitions. The interface supports both free exploration and scripted scenario playback for expert analysis and public demonstrations alike. Figure~\ref{fig:structure} illustrates the overall architecture.

\begin{figure}[t]
    \centering
    \includegraphics[width=0.95\linewidth]{co-sim.pdf}
    \caption{System architecture: SUMO and PyPSA synchronized by event-driven coordination with progressive time and vehicle quota manager. LLM agent (GPT-4) interprets natural language and controls Mapbox visualization.}
    \label{fig:structure}
\end{figure}

\section{Demonstration Scenarios}
We illustrate the system through three scenarios highlighting power-transportation interdependence. Scenarios are launched via natural language commands with effects visualized in real-time (Figure~\ref{fig:interface}).

\begin{figure}[t]
    \centering
    \includegraphics[width=0.95\linewidth]{Screenshot.png}
    \caption{LLM-to-Map interface showing Manhattan testbed with substations, traffic signals, and conversational chatbot control panel.}
    \label{fig:interface}
\end{figure}

\noindent\textbf{Scenario 1: Grid-Induced Charging Disruptions.}
A user types "turn off Penn Station substation," immediately disabling all charging stations in that service area. Vehicles attempting to charge reroute to operational stations, creating localized queues and circling behavior. Vehicles unable to secure charging risk battery depletion and stranding. PyPSA reflects load loss and demand redistribution. Users can restore substations via "restore Penn Station" or "restore all" commands, instantly re-enabling charging infrastructure. This scenario demonstrates how single grid outages cascade into widespread mobility disruptions (Figure~\ref{fig:scenario1}).

\begin{figure}[t]
\centering
\includegraphics[width=0.9\columnwidth]{scenario1.png}
\caption{Scenario 1: Grid-induced charging disruption showing (1) normal operations and (2) Penn Station offline after command execution.}
\label{fig:scenario1}
\end{figure}

\noindent\textbf{Scenario 2: EV Charging Stress.}
This scenario explores how high EV penetration stresses distribution substations. Vehicles spawn with 25–30\% battery, immediately seeking charging. Users trigger "blackout scenario," failing all substations except one. The camera zooms out to show sequential failures across the network, then focuses on the sole operational station. As simulation time progresses (e.g., 08:00 $\rightarrow$ 08:30), the progressive time system automatically maintains rush-hour vehicle quotas (85–100), continuously spawning replacements as vehicles complete trips. This sustained pressure persists throughout the rush period, with all EVs converging on one location. The station rapidly reaches capacity, creating visible queuing patterns with vehicles circling while the camera captures their desperate search for charging. Some become stranded as batteries deplete. This demonstrates the feedback loop from mobility behavior into grid stress, emphasizing infrastructure redundancy importance during emergencies (Figure~\ref{fig:scenario2}).

\begin{figure}[t]
\centering
\includegraphics[width=0.9\columnwidth]{scenario2.png}
\caption{Scenario 2: Blackout stress test showing (1) normal grid, (2) 7/8 substations failed, (3) EVs converging on single station, and (4) vehicles stranded.}
\label{fig:scenario2}
\end{figure}

\noindent\textbf{Scenario 3: Vehicle-to-Grid Rescue.}
V2G is activated as a resilience measure. Users type "run V2G scenario," triggering Times Square substation failure with people trapped in an elevator. The system identifies nearby EVs with battery $\geq$70\%, sending compensation-incentivized discharge requests. Vehicles spawn with high battery levels to ensure participation. The camera zooms to the failed substation, showing EVs responding, connecting to discharge points, and collectively injecting power. Aggregated injection raises voltages, restores substation operation, frees trapped occupants, and brings signals online. Mapbox visualizations show affected EVs rerouting to V2G-enabled stations and discharging energy, with restored signals flashing back to green as power returns. The chatbot provides comprehensive session summaries including total power restored, participating vehicles, and compensation earned. This scenario demonstrates how EVs transform from grid stressors to stabilizers during emergencies (Figure~\ref{fig:scenario3}).

\begin{figure}[t]
\centering
\includegraphics[width=0.9\columnwidth]{scenario3.png}
\caption{Scenario 3: V2G emergency response showing (1) normal state, (2) substation failure with trapped people, (3) EVs discharging to grid, and (4) power restored.}
\label{fig:scenario3}
\end{figure}

\noindent\textbf{Interactive Exploration.}
While scenarios can be played in scripted form, the demo supports exploration. Attendees may combine scenarios or issue free-form queries such as "zoom to Grand Central and show traffic impact." The LLM agent interprets conversational inputs, executes simulator actions, and returns synchronized visualizations. This transforms the demo into a live laboratory for studying cascading failures and resilience strategies.

\section{Discussion and Outlook}
The \textbf{LLM-to-Map} system demonstrates how conversational interfaces bridge the gap between technical co-simulation tools and accessible urban analysis. By integrating SUMO, PyPSA, and Mapbox under a common event-driven engine with progressive temporal dynamics, the system captures bidirectional effects between networks while maintaining realistic traffic patterns through intelligent resource management. The LLM agent lowers interaction barriers, enabling users to trigger outages, observe cascading failures, and activate V2G responses without domain expertise.

The architecture is generalizable: any city with OpenStreetMap and utility data can be incorporated, and the modular design allows additional infrastructures (communication, water, supply chains) to be added as plugins. From a research perspective, the system contributes to both web search (fuzzy geospatial retrieval, conversational state tracking) and urban computing (interactive resilience analysis). For practitioners, it offers hands-on decision support making cascading risks and mitigation strategies visible in real time. We envision LLM-to-Map evolving into a collaborative platform where researchers, operators, and citizens explore smart city resilience together.

\bibliographystyle{ACM-Reference-Format}
\bibliography{z}

\end{document}
\endinput
%%
%% End of file `sample-acmcp.tex'.
